% OAI O-RAN 7.2 Lab -- End-to-End Debugging Technical Report
% Compile with pdfLaTeX on Overleaf.
\documentclass[11pt,a4paper]{article}

\usepackage[T1]{fontenc}
\usepackage[utf8]{inputenc}
\usepackage[margin=1in]{geometry}
\usepackage{booktabs}
\usepackage{tabularx}
\usepackage{array}
\usepackage{amssymb,amsmath}
\usepackage{xcolor}
\usepackage{listings}
\usepackage{enumitem}
\usepackage{caption}
\usepackage[hidelinks]{hyperref}
\usepackage{underscore}   % after hyperref: _ prints literally in text/\texttt

% ---- colour macros ----
\definecolor{closedcol}{rgb}{0.05,0.45,0.05}
\definecolor{opencol}{rgb}{0.70,0.15,0.10}
\definecolor{codebg}{gray}{0.96}

\newcommand{\closed}{\textcolor{closedcol}{\textbf{[Closed]}}}
\newcommand{\openissue}{\textcolor{opencol}{\textbf{[Open]}}}
\newcommand{\yes}{\textcolor{closedcol}{\ensuremath{\checkmark}}}
\newcommand{\warn}{\textcolor{orange!85!black}{\textbf{(!)}}}
\newcommand{\no}{\textcolor{opencol}{\ensuremath{\times}}}
\newcommand{\us}{\ensuremath{\mu}\text{s}}

% ---- listing style ----
\lstdefinestyle{plain}{
  basicstyle=\ttfamily\small,
  backgroundcolor=\color{codebg},
  breaklines=true,
  columns=fullflexible,
  keepspaces=true,
  frame=single,
  framesep=4pt,
  xleftmargin=2pt,
  showstringspaces=false,
}
\lstset{style=plain}

\setlength{\parskip}{0.45em}
\setlength{\parindent}{0pt}
\renewcommand{\arraystretch}{1.18}

\title{\textbf{OAI O-RAN 7.2 Lab\\[4pt]End-to-End Debugging Technical Report}}
\author{Jesse}
\date{2026-05-26 $\rightarrow$ 2026-06-01}

\begin{document}
\maketitle

\paragraph{Summary.}
An OAI 5G NR O-RAN 7.2 split cell (DU + O-RU + UE, single host, DPDK/SR-IOV fronthaul,
vrtsim air channel) was brought up from zero through thirteen tracked issues. Experiments
swept IQ compression width (iq8--16), MIMO count (1$\times$1--4$\times$4), receive SNR,
DL fronthaul latency (T2a\_min), and UE mobility. Key results: the 10~MHz cell delivers
$\sim$5~Mbps UL at 0\% loss; FH load scales linearly with compression but UL is flat
(fronthaul over-provisioned 14$\times$); SNR degrades UL monotonically (20~dB
$\rightarrow$ 5~Mbps, 8~dB $\rightarrow$ 2.1~Mbps); tightening the DL FH window shows a
sharp knee at 400--450~\us; all attempts beyond 10~MHz fail.

\tableofcontents
\bigskip\hrule\bigskip

%% ==========================================================================
\section{Overview and Architecture}
%% ==========================================================================

\subsection{System diagram}

\begin{figure}[h]
\centering
\scalebox{0.80}{
\begin{minipage}{1.25\linewidth}
\begin{lstlisting}[breaklines=false,basicstyle=\ttfamily\footnotesize,frame=single,
                   caption={},belowskip=0pt,aboveskip=0pt]
  +-------------+  O-RAN 7.2 FH (eCPRI/xran     +---------+  vrtsim "air"  +------------------+
  | DU          |  DPDK SR-IOV, 2 VLANs)         | O-RU    |  shared mem.   | UE               |
  | nr-softmodem|<==============================>| nr-oru  |<=============> | nr-uesoftmodem   |
  | (PHY/MAC)   |  DL:vf2/vf3   UL:vf0/vf1      | (RUPHY) |  chanmod/IQ    | oaitun_ue1       |
  +------+------+  10G VEB, PF link DOWN          +---------+                +--------+---------+
         |                                                                            |
         | NGAP/GTP-U                                              iperf3 -c 10.0.0.1 -u -b 20M
         v                                                                            |
  +------------------+                                                                |
  | OAI 5GC          |  docker-compose: AMF/SMF/NRF/UPF                              |
  | UPF  10.0.0.1    |<--------------------------------------------------------------+
  | iperf3 -s        |  GTP-U decap          TRUE UL = server received rate (~5 Mbps)
  | (UPF netns)      |
  +------------------+
\end{lstlisting}
\end{minipage}
}
\caption{System architecture with iperf3 endpoints. UL goodput = server-side received rate.}
\label{fig:arch}
\end{figure}

\subsection{Two interfaces --- keep them separate}

\begin{table}[h]\centering\small
\begin{tabularx}{\linewidth}{@{}l l X@{}}
\toprule
\textbf{Interface} & \textbf{Medium} & \textbf{Bug class found here} \\
\midrule
Fronthaul (FH) DU$\leftrightarrow$RU & O-RAN 7.2 eCPRI/xran, DPDK SR-IOV &
  C-plane timing, slot-ID mismatch, U-plane buffer reset, compression header \\
Air RU$\leftrightarrow$UE & vrtsim shared memory &
  chanmod realtime, SNR, mobility \\
\bottomrule
\end{tabularx}
\end{table}

\subsection{Cell and hardware}

\begin{itemize}[leftmargin=1.4em]
\item \textbf{Cell:} 24~PRB, $\mu$1 (30~kHz SCS) $\Rightarrow$ \textbf{10~MHz}, fs = 15.36~Msps,
  band n77, fc = 4049.76~MHz.
\item \textbf{NIC:} Intel X710 (i40e), PF \texttt{eno1np0} (\texttt{0000:05:00.0}), PF link
  DOWN --- loopback via embedded switch (VEB), no cable.
  VF map: vf0/vf1 = RU egress (UL U-plane, VLAN 3/4);
  vf2/vf3 = DU egress (DL U-plane, VLAN 3/4).
\item \textbf{Build trap:} the FH code (\texttt{oaioran.c}, \texttt{oaioran_ru.c},
  \texttt{oai_bfp_compression.c}, \texttt{oran-config.c}) compiles into a
  runtime-dlopen'd shared library \texttt{liboran_fhlib_5g.so} --- \emph{not} into the
  \texttt{nr-softmodem} or \texttt{nr-oru} executables. Rebuilding the executables does not
  pick up FH changes; \texttt{ninja oran_fhlib_5g} is required. Several ``fix had no effect''
  episodes traced back to this.
\end{itemize}

%% ==========================================================================
\section{Issues}
%% ==========================================================================

Each subsection is labelled \closed\ (fixed, in the working build) or
\openissue\ (known limitation, no fix yet merged).

%% --------------------------------------------------------------------------
\subsection{\texorpdfstring{\closed}{[Closed]}\ AVX512 unavailable --- BFP compression reimplemented without AVX512}
%% --------------------------------------------------------------------------

\textbf{Symptom.}
The initial BFP (Block Floating Point) compression implementation relied on AVX-512 SIMD
intrinsics. The lab host does not expose AVX-512 (either unsupported or not enabled in the
BIOS/kernel). This caused the BFP code path to fail --- either as an illegal-instruction
fault at runtime or as silent miscompilation where the intrinsics were never executed ---
making BFP compression non-functional from the outset.

\textbf{Root cause.}
The compression library (\texttt{oai_bfp_compression.c}) was written with AVX-512 as the
assumed baseline SIMD tier. On hardware that only exposes up to AVX2 (or on builds where the
compiler does not emit AVX-512), the data path either trapped or produced garbage IQ output.

\textbf{Fix (commit \texttt{40eaf33c9c}).}
BFP compression was reimplemented using portable scalar code, keeping the same
O-RAN-compliant bit-packing format. The scalar path is correct and produces identical output
to the AVX-512 path. This is the baseline that all subsequent compression work (the
compression-header fix in Section~2.6) was built on.

%% --------------------------------------------------------------------------
\subsection{\texorpdfstring{\closed}{[Closed]}\ DU/RU binary mismatch --- invalid U-plane slot IDs}
%% --------------------------------------------------------------------------

\textbf{Symptom.}
\begin{lstlisting}
[ORU] Drop invalid U-plane packet: ... slot=38 ... symbol=37 num_prb=0
Packets early: 1  /  Packets malformed: 1
\end{lstlisting}
For $\mu$1, valid slot\,=\,0--19 and symbol\,=\,0--13.
\texttt{slot=38} and \texttt{symbol=37} are impossible from a correct DU.

\textbf{Root cause.}
\texttt{run_ru.sh} launched the new \texttt{build/nr-oru}, while \texttt{run_du.sh} still
launched the old \texttt{cmake_targets/ran_build/build/nr-softmodem}.
\texttt{xran_slotid_convert()} in \texttt{phy-f-1.0/.../xran_main.c} is a no-op
(\texttt{return slot_id;}) in the current source, but the old DU binary was linked against a
build where the $\mu$1 shift (\texttt{slot\_id <{}< 1}) was active. The DU encoded wire slot
$19 << 1 = 38$; the RU applied the no-op and read 38 --- out of bounds --- and dropped the
packet. Confirmed not to be a stray ARP/IPv6 frame: xran's \texttt{handle_ecpri_ethertype()}
only dispatches EtherType \texttt{0xAEFE}, so non-eCPRI frames are discarded before the
packet processor; the offending packet passed the full eCPRI header parse.

\textbf{Fix.}
Point both launchers at \texttt{openairinterface5g/build/} and rebuild DU + RU from the same
\texttt{phy-f-1.0} source tree so \texttt{xran_slotid_convert} is the no-op in both
binaries. Drop counter (\texttt{up_malformed}/\texttt{up_dropped}, bounds-checked at
\texttt{oaioran_ru.c:640}) went to 0.

%% --------------------------------------------------------------------------
\subsection{\texorpdfstring{\closed}{[Closed]}\ PRACH MSG1 not forwarded --- MSG2 (RAR) never sent}
\label{sec:prach}
%% --------------------------------------------------------------------------

\textbf{Symptom.}
After DL synchronization, the UE repeatedly transmitted PRACH (MSG1) but received no RAR
(MSG2). The gNB never logged a PRACH detection event; the DU never scheduled a RAR because
it never received the complete preamble IQ from the RU. The UE exhausted its preamble
transmission counter and re-entered idle.

Three independent bugs on the FH C-plane path caused the RU to silently discard every PRACH
IQ packet before forwarding it to the DU.

\paragraph{(A) \texttt{T1a\_cp\_ul} timing too tight.}
The DU must send a C-plane section-type-3 packet to the RU \emph{before} each PRACH slot; the
RU caches it in \texttt{prach_config_by_frame_slot[...]} and consults it in
\texttt{xran_oru_send_prach()}. With \texttt{T1a_max_cp_ul = 429}\,\us:
\begin{lstlisting}
delay_cp_ul = 500 - 429 = 71 us  ->  C-plane fired at symbol 2 of the PRACH slot
\end{lstlisting}
The RU's read thread processed symbols 0 and 1 with an empty cache --- those PRACH symbols
were silently dropped. PRACH format~B4 (\texttt{prach_config_index=159}) needs \textbf{12
consecutive symbols}; losing symbols 0--1 breaks preamble detection entirely.

\textbf{Fix:} \texttt{T1a_cp_ul = (285, 535)} in \texttt{du_test.conf}. With
\texttt{T1a\_max > 500}\,\us, xran's \texttt{xran_timing_create_cbs()} wraps one slot
earlier (\texttt{offset_num_slots_cp_ul = 1}) and fires the C-plane at symbol~0 of slot~18,
populating the cache $\sim$450\,\us\ before slot~19 begins. Config-only, no recompile.

\paragraph{(B) C-plane packet built but never transmitted.}
\texttt{tx_cp_ul_cb()} (\texttt{xran_main.c:1401}) checked:
\begin{lstlisting}[language=C]
if (ret == XRAN_STATUS_SUCCESS)   // XRAN_STATUS_SUCCESS == 0
    send_cpmsg(...);
\end{lstlisting}
but \texttt{generate_cpmsg_prach()} returns a \textbf{positive byte count} on success, so
\texttt{send_cpmsg()} was never reached --- the packet was assembled and the mbuf allocated,
then discarded without being placed in the TX ring.

\textbf{Fix:} \texttt{if (ret >= 0)}.

\paragraph{(C) Slot index overflow in the RU PRACH cache write.}
\texttt{process_ru_cplane()} (\texttt{oaioran_ru.c:700}) derived the numerology from
\texttt{frameStructure.uScs}, which for format~B4 encodes the \textbf{PRACH SCS (= 14)}, not
$\mu$:
\begin{lstlisting}[language=C]
int mu   = hdr->frameStructure.uScs;      // 14, not 1
int slot = slotId + subframeId*(1<<mu);   // 1 + 9*16384 = 147457  -> out of bounds
\end{lstlisting}
The cache entry was never written, so \texttt{xran_oru_send_prach()} returned early on every
PRACH symbol.

\textbf{Fix:} \texttt{int mu = fh\_cfg->frame\_conf.nNumerology;}

Also fixed during this period: \texttt{callbacks_per_slot} in \texttt{oran-init.c} was
\texttt{2}, raised to \textbf{14} (one RU consumer wakeup per OFDM symbol).

After (A) + (B) + (C): PRACH IQ forwarded intact, the gNB detected the preamble,
scheduled a RAR, and MSG2 was transmitted to the UE. Rebuild scope:
\texttt{libxran.so} for (A,B); \texttt{liboran_fhlib_5g.so} for (C).

%% --------------------------------------------------------------------------
\subsection{\texorpdfstring{\closed}{[Closed]}\ MSG3 contention resolution failure --- UL U-plane buffer reset wiped early symbols}
\label{sec:ulbuf}
%% --------------------------------------------------------------------------

\textbf{Symptom.}
After the PRACH fixes the UE received MSG2 (RAR) with a TC-RNTI and a UL grant for MSG3.
The UE transmitted MSG3 (RRC Setup Request) on the scheduled UL PUSCH resource, but the gNB
could not decode it. MSG4 (Contention Resolution) was never sent, the MAC contention
resolution timer expired, and the UE re-initiated RA from MSG1 in a continuous loop.
Even after RA appeared to eventually complete (NAS attach), UL goodput was effectively zero:
UE MAC counter \texttt{UL harq: \textasciitilde719/718} = 99.9\% HARQ retransmissions at
$\sim$34~dB SNR. MSG3 was simply the first symptom of the same underlying defect affecting
all UL PUSCH transmissions.

\textbf{Root cause.}
The RU streams UL U-plane symbol-by-symbol in near real time. An uncommitted line in
\texttt{oai_xran_fh_rx_callback} (\texttt{radio/fhi_72/oaioran.c}) reset the \emph{next}
slot's ring-buffer entry, \texttt{(tti + 1) \% 20}, from the mid-slot callback
(\texttt{rx\_sym == 7}). But the next slot's first U-plane packets \emph{arrive before that
callback fires}: by the time the mid-slot fires and clears \texttt{(tti+1)}'s
\texttt{nSecDesc}, symbols 0--8 had already landed in that buffer and were wiped. Only the
last $\sim$5 symbols of each full UL slot survived into the DU's PUSCH decoder.

\textbf{Evidence.} A temporary \texttt{[UL SYMPROBE]} in \texttt{xran_fh_rx_read_slot}
logged \texttt{present=5/14 mask=00000000011111} (syms 9--13 only) --- deterministic, not
random. A \texttt{[FH DROP]} counter confirmed xran was not dropping
(\texttt{Rx\_pkt\_dupl=0}, \texttt{Rx\_on\_time} large) --- the symbols reached the buffer;
the reset erased them. The LDPC decoder corroborated: UL TBs decoded only on HARQ round~3
(rv3), with \texttt{ulsch\_llr[0..3]=0,0,0,0} (the missing early-symbol LLRs defaulting to
zero); REs accumulated across three HARQ rounds before a codeword decoded, which for MSG3
(a short one-shot TB) meant the contention resolution window always expired first.

\textbf{Fix (commit \texttt{37b55fa535}).}
\begin{lstlisting}[language=C]
// oaioran.c  oai_xran_fh_rx_callback
- reset_buffer((tti + 1) % 20);
+ reset_buffer((tti + XRAN_N_FE_BUF_LEN/2) % 20);   // tti+10, half a ring away
\end{lstlisting}
\texttt{tti+10} is cleared long before its slot is filled or read, and never overlaps the
active fill/read window.

\textbf{Result:} \texttt{present=14/14 mask=11111111111111}, all HARQ round~0,
\texttt{UL harq: 845/0} --- \textbf{zero retransmissions}. MSG3 decoded on the first
transmission; contention resolution completed; NAS registration (Auth / SecurityMode /
Registration Accept) now completed end-to-end.

%% --------------------------------------------------------------------------
\subsection{\texorpdfstring{\closed}{[Closed]}\ iq16 path unstable --- PUSCH compression-header byte offset}
%% --------------------------------------------------------------------------

\textbf{Symptom.}
With UL flowing, iq9 (BFP-compressed) was stable; iq16 (uncompressed,
\texttt{compMeth = NONE}) showed wildly varying UL SNR: $-8.5$ to $+39$~dB vs.\ a stable
63~dB for iq9.

\textbf{Root cause.}
\texttt{xran_oru_send_pusch()} (\texttt{oaioran_ru.c}) unconditionally wrote a 2-byte
\texttt{data_section_compression_hdr} before the IQ payload, even for \texttt{compMeth =
NONE}. The DU's xran reader (\texttt{xran_common.c:322}) computes
$\texttt{iq\_offset} = \texttt{ecpri} + \texttt{radio\_app} + \texttt{data\_section}$ for
NONE (no compression header), so it read those 2 zero bytes as the \textbf{first IQ sample},
shifting every UL sample by one position. For iq9/BFP, \texttt{xran_common.c} adds
\texttt{compr\_size} to \texttt{iq\_offset}, so the alignment was always correct --- which is
why only iq16 failed.

\textbf{Fix (commit \texttt{47dcaef44b}).}
Gate the header write: \texttt{use\_comp\_hdr = (compMeth != NONE)}.
Also fixed: a BFP mbuf under-allocation (2 bytes short, truncating the tail of each BFP
packet). iq16's truncation had been masked in iq9 runs by
\texttt{ul\_prbblacklist = "0,1,2,3,20,21,22,23"}, which hides the edge PRBs where the
truncation fell.

%% --------------------------------------------------------------------------
\subsection{\texorpdfstring{\closed}{[Closed]}\ PDU session reject --- SMF deadlock in CN}
%% --------------------------------------------------------------------------

\textbf{Symptom.}
UE attached at RRC but received \texttt{FGS\_PDU\_SESSION\_ESTABLISHMENT\_REJ}; no IP
address, no \texttt{oaitun\_ue1}.

\textbf{Root cause.}
The SMF had deadlocked $\sim$31~h earlier (last log 2026-05-29 03:40; likely blocked on a PFCP
recv from the UPF). The container reported healthy but the application was frozen, emitting
no NRF heartbeats. Consequently the NRF had no SMF registered
(\texttt{/nnrf-nfm/v1/nf-instances?nf-type=SMF} $\rightarrow$ \texttt{item:[]}). The AMF
(\texttt{enable\_smf\_selection: yes}) could not locate an SMF and logged \emph{``No SMF is
available for this PDU session''}.

\textbf{Fix.}
Full CN restart --- a bare SMF restart fixed NRF discovery but the first session still
rejected with ``empty QFI list'' from stale UPF/N4 state; the full down/up cleared it:
\begin{lstlisting}[language=bash]
cd ~/oran_lab/oai-cn5g && docker-compose down && docker-compose up -d   # v1 syntax
\end{lstlisting}
DNN ``oai'', sst=1, 5QI~9, PAA pool 10.0.0.0/24 were all correct. After restart: UE gets
10.0.0.x, TUN up, traffic flows.

\textbf{End-to-end milestone (2026-05-30):} iq9 UL 1.325~Mbps / DL 6.082~Mbps;
iq16 UL 1.325~Mbps / DL 7.130~Mbps.

%% --------------------------------------------------------------------------
\subsection{\texorpdfstring{\closed}{[Closed]}\ UL grant pinned at 5 RBs --- \texttt{min\_grant\_prb} default}
%% --------------------------------------------------------------------------

\textbf{Symptom.}
After all radio/CN fixes, UL throughput was stuck at $\sim$1.3~Mbps.
The UE MAC showed grants of only 5~RBs regardless of how much data was queued.

\textbf{Root cause.}
\texttt{min\_grant\_prb} in MACRLCs defaults to \textbf{5}
(\texttt{MACRLC\_nr\_paramdef.h}, \texttt{.defintval=5}). The UE is perpetually scheduled
\texttt{sched\_inactive}: its transmit buffer drains every grant, so the scheduler's estimate
\texttt{B = estimated\_ul\_buffer - sched\_ul\_bytes == 0}
(\texttt{gNB\_scheduler\_ulsch.c:2113}), and inactive UEs always receive exactly
\texttt{min\_grant\_prb} RBs (lines $\sim$2190/2260). A ``greedy-RB'' hack in the
\texttt{B>0} branch was tried but had no effect --- that branch is never entered when the
buffer drains on every TTI.

\textbf{Fix.}
\texttt{min\_grant\_prb = 16;} in \texttt{du\_test.conf} MACRLCs.
UL grant jumped 5 $\rightarrow$ 17~RBs.

%% --------------------------------------------------------------------------
\subsection{\texorpdfstring{\closed}{[Closed]}\ DL-heavy TDD starved UL --- pattern changed to D\,M\,U\,U\,U}
%% --------------------------------------------------------------------------

The default TDD pattern \texttt{D D D M U} gives only 1 UL slot per 5-slot period.
Changed to \textbf{\texttt{D M U U U}} in both \texttt{du\_test.conf}
(\texttt{nrofDownlinkSlots 3$\rightarrow$1}, \texttt{nrofUplinkSlots 1$\rightarrow$3}) and
\texttt{ru\_test.conf} (\texttt{num\_dl\_slots 3$\rightarrow$1},
\texttt{num\_ul\_slots 1$\rightarrow$3}). The UE follows via SIB1.
UL TBs per run went $\sim$2337 $\rightarrow$ 26\,000+.

Combined with the \texttt{min\_grant\_prb} fix: UL $\approx$\textbf{5.05~Mbps @ 0\% loss,
17~RBs, 0 retx} (from 1.3~Mbps, a $\sim$3.9$\times$ improvement).

\textbf{Note --- TCP iperf breaks with this pattern.} DL is so starved that TCP ACKs and
the iperf control socket reset (``Connection reset by peer''). UL \emph{must} be measured
with UDP (see Section~\ref{sec:iperf}).

%% --------------------------------------------------------------------------
\subsection{\texorpdfstring{\closed}{[Closed]}\ iperf measurement traps --- how server-side UL is measured}
\label{sec:iperf}
%% --------------------------------------------------------------------------

Getting a trustworthy UL figure required resolving five measurement pitfalls, each of which
produced an incorrect reading at some point.

\textbf{Pitfall 1 --- wrong server address inflates by $\sim$7000$\times$.}
Pointing the iperf3 client at the host IP (\texttt{192.168.70.129}) instead of the UPF
gateway (\texttt{10.0.0.1}) causes iperf to loop back on the host without crossing the
radio, reporting a bogus $\sim$36~Gbps. Traffic to \texttt{10.0.0.1} is forced over the radio
because the host routes \texttt{10.0.0.0/24} via \texttt{oaitun\_ue1}. The server must bind
\texttt{10.0.0.1}.

\textbf{Pitfall 2 --- UPF container has no iperf3.}
The host's iperf3 is run inside the UPF network namespace (giving it ownership of
\texttt{tun0}/\texttt{10.0.0.1}) while keeping the host mount namespace
(\texttt{/tmp} stays on the host):
\begin{lstlisting}[language=bash]
UPF_PID=$(docker inspect -f '{{.State.Pid}}' oai-upf)
sudo nsenter -t $UPF_PID -n pkill -x iperf3           # clear stale servers
sudo nsenter -t $UPF_PID -n iperf3 -s -B 10.0.0.1 \
    -1 --logfile /tmp/iperf_srv.log -D                 # one-shot, daemonized
iperf3 -c 10.0.0.1 -u -b 20M -t 20                   # UE side (client)
\end{lstlisting}

\textbf{Pitfall 3 --- UDP client rate is not the radio rate.}
The client's \texttt{bits\_per\_second} over UDP is the rate \emph{pushed into the TUN}; the
UE drops excess at the TUN when the radio can't keep up, so the client over-reports (source
of the fake ``20~Mbps'' readings). The true UL rate is the \textbf{server's received rate},
read from \texttt{/tmp/iperf\_srv.log}. The sweep uses \texttt{IPERF\_UDP=1
IPERF\_UDP\_RATE=20M} and parses the server log.

\textbf{Pitfall 4 --- \texttt{end.sum} is empty under saturation.}
When the UL saturates the radio, the iperf3 control socket dies before the JSON summary is
written. \texttt{--get-server-output}/\texttt{end.sum.bits\_per\_second} come back blank.
The harness instead averages the per-interval rows for throughput, jitter, and loss.

\textbf{Pitfall 5 --- two concurrent sweeps collide.}
The harness runs a single iperf3 server (\texttt{ensure\_iperf\_server} kills duplicates and
starts one). A second sweep returns ``server is busy'' $\rightarrow$ 0~Mbps. The two sweeps
also share \texttt{ru\_test.conf}/\texttt{du\_test.conf} and vrtsim cores. Never run two
sweeps at the same time.

\textbf{Cross-check.} The UE MAC counter \texttt{UL harq: X/Y} $\times$ TBS
($\sim$1.5~KB/TB at 17~RB / 64-QAM) independently gives $\sim$5~Mbps, confirming the iperf
server number.

%% --------------------------------------------------------------------------
\subsection{\texorpdfstring{\openissue}{[Open]}\ 10 MHz bandwidth ceiling --- three blocking layers}
%% --------------------------------------------------------------------------

The cell is 24~PRB (10~MHz) by construction. Every path to wider bandwidth is blocked:

\textbf{Layer 1 --- PRB blacklist (config, effect is immediate).}
\texttt{ul\_prbblacklist = "0,1,2,3,20,21,22,23"} hides 8 of 24~PRBs, leaving only
16~usable, and MCS is already maxed at 28. So the UL ceiling is
$16\,\text{PRB} \times \text{MCS28} \times \text{UL-TDD duty cycle} \approx 5\,\text{Mbps}$
within 10~MHz. The blacklist exists because the edge PRBs were where the BFP truncation and
band-edge artefacts landed (Section~2.5).

\textbf{Layer 2 --- FH-tuned PHY breaks SIB1 decode outside 24~PRB.}
The SSB/pointA frequency placement was fixed in config
(\texttt{dl\_absoluteFrequencyPointA = 669984 - (BW\_PRB/2)*24},
UE \texttt{--ssb = (BW\_PRB/2 - 10)*12}), and at 106~PRB the UE now syncs and finds the
CORESET0/SIB1 grant. But SIB1 PDSCH still NACKs
(\texttt{Got NACK on NR-BCCH-DL-SCH-Message (SIB1)}): the lab's FH-tuned PHY commit
\texttt{3b10464162} (amplitude/timing tuned for the 7.2 freq-domain split where the \emph{RU}
does the IFFT) produces confident-but-wrong LLRs outside 24~PRB
(\texttt{l\_max=127} saturated, 8/8 LDPC iterations, wrong codeword). The rfsim path
exhibits the same failure (Section~2.13).

\textbf{Layer 3 --- realtime constraints.}
At 106~PRB realtime vrtsim+xran, the gNB runs N\_RB~106 and SIB1 decodes clean, but
\texttt{prach\_I0 = 0.0~dB} at the gNB --- RA never completes. The scalar chanmod FIR
(\texttt{vrtsim.c:887}) also cannot sustain 4--8$\times$ the 10-MHz sample rate without
overrunning its core (Section~2.12). At 273~PRB (100~MHz) the DU crashes with a realtime
timeout.

\textbf{FH headroom is not the limit.} Direct testpmd VF$\leftrightarrow$VF measurement:
$\sim$3.18~Gbps per VF, $\sim$6.36~Gbps aggregate, zero loss. OAI loads the FH to
$\sim$440~Mbps ($<$15\%). FH lead-time is flat at $\sim$12 symbols ($\sim$428\,\us\ margin)
across iq8--iq16 (94--180~Mbps). The FH pipe is not the bottleneck; the per-slot timing
window is.

%% --------------------------------------------------------------------------
\subsection{\texorpdfstring{\openissue}{[Open]}\ $\sim$50\% attach flakiness --- per-bringup timing transient}
%% --------------------------------------------------------------------------

Across $\sim$250 logged runs, roughly half end in \texttt{ue-timeout}: either the UE never
achieves initial DL sync ($\sim$23\% of runs, 0--200 \texttt{synch Failed} in UE log before
giving up) or it syncs but PRACH fails (\texttt{prach\_I0 = 0.0~dB} at the gNB, $\sim$27\%
of runs). Once the UE locks it stays locked for the whole session --- the failure is a
startup/timing-alignment transient, not a steady-state instability. The experiment harness
mitigates with a retry-to-attach loop; all results in Section~4 are from successful
bringups only.

%% --------------------------------------------------------------------------
\subsection{\texorpdfstring{\openissue}{[Open]}\ Scalar chanmod FIR limits 4$\times$4 and wide-BW realtime}
%% --------------------------------------------------------------------------

The channel-modelling convolution \texttt{perform\_channel\_modelling()} in
\texttt{radio/vrtsim/vrtsim.c:887} is a scalar, un-vectorised complex FIR loop (comment:
\texttt{// TODO: Use AVX2 for this}). Cost scales as $N_\text{tx} \times N_\text{samp} \times
N_\text{taps}$ per actor thread per slot. At 1$\times$1 TDL-D it is borderline on one core;
at 4$\times$4 the server-side DL chanmod runs 16 convolutions/slot and consistently overruns
(``Packets not processed by the application layer (application layer too slow)''), causing the
gNB to read zero-filled UL buffers ($\Rightarrow$ \texttt{prach\_I0=0.0~dB} $\Rightarrow$ no
RA). A core-pin fix was applied (pin chanmod actors to isolated cores 24+/28+ instead of
\texttt{init\_actor(...,-1)}, which let them float onto busy cores 0--3) --- this unblocked
1$\times$1 mobility and 2$\times$2, but 4$\times$4 / wide-BW still overrun. The fix required
is to vectorise the FIR with AVX2/AVX512.

%% --------------------------------------------------------------------------
\subsection{\texorpdfstring{\openissue}{[Open]}\ rfsim DL broken on FH-tuned tree}
%% --------------------------------------------------------------------------

To isolate PHY bugs from FH bugs, the rfsim path (no FH, monolithic gNB) was tested.
rfsim DL fails: SIB1 LDPC produces confident-but-wrong LLRs (same signature as the wide-BW
failure in Section~2.10). vrtsim is immune because the split RU does the IFFT on a separate
path; rfsim uses the monolithic gNB IFFT, which is affected by a suspected
\texttt{txdataF} fftshift/flat-buffer refactor (\texttt{044ca835d6}) that interacts badly
with the FH-tuned amplitude/timing commit \texttt{3b10464162}.

A clean OAI tree (\texttt{git worktree} at release \texttt{2026.w22}, commit
\texttt{26efcc4989}, where rfsim DL is a hard CI gate) was checked out at
\texttt{oai\_rfsim\_clean/} for high-BW / non-realtime work, keeping the FH-tuned tree
for the FH measurements. The two halves are independent: rfsim needs no xran; FH latency
can be injected synthetically from a NIC-measured load$\rightarrow$latency LUT.

%% ==========================================================================
\section{Toolchain}
%% ==========================================================================

The lab is orchestrated by a set of shell scripts in \texttt{oran\_lab/}.
All scripts must be run as the regular user (not \texttt{sudo}) --- they invoke \texttt{sudo}
internally for DPDK and network configuration. Running as root breaks path resolution because
\texttt{HOME=/root} redirects away from the user's environment.

%% --------------------------------------------------------------------------
\subsection{\texttt{prepare\_network.sh} --- DPDK and VF setup (run once at boot)}
%% --------------------------------------------------------------------------

Prepares the SR-IOV and hugepage environment before any OAI process is started:
\begin{itemize}[leftmargin=1.4em]
\item Creates 5 VFs on PF \texttt{eno1np0}, sets their MAC addresses and VLAN tags
  (vf0: MAC \texttt{...66}, VLAN 3; vf1: \texttt{...67}, VLAN 4; etc.).
\item Binds each VF to the \texttt{uio\_pci\_generic} driver so DPDK can take ownership.
\item Allocates 8192 hugepages (2~MB each, 16~GB total) and mounts
  \texttt{/dev/hugepages}.
\item Configures the embedded switch so vf0/vf1 (RU egress) can reach vf2/vf3 (DU egress)
  across the VEB --- the PF link is intentionally left DOWN; no physical cable is needed.
\end{itemize}

Between runs, stale DPDK hugepage mappings must be cleared manually:
\texttt{sudo rm -f /dev/hugepages/*map\_*} (DPDK does not release them on process kill).

%% --------------------------------------------------------------------------
\subsection{\texttt{run\_du.sh} --- launch the DU (nr-softmodem)}
%% --------------------------------------------------------------------------

Launches \texttt{openairinterface5g/build/nr-softmodem} with:
\begin{itemize}[leftmargin=1.4em]
\item \texttt{--sa} (standalone 5G), \texttt{--rfsim} disabled, fronthaul via
  \texttt{liboran\_fhlib\_5g.so} (loaded at runtime).
\item Config file \texttt{du\_test.conf} (MACRLCs, cell parameters, FH timing windows
  \texttt{T1a\_cp\_ul}, \texttt{T2a\_up}, TDD pattern, \texttt{min\_grant\_prb}, etc.).
\item CPU affinity to isolated cores 4--9, 17, 18 (away from the OS and from the RU/UE
  process sets).
\item Log file \texttt{logs/run\_du/<timestamp>/du.log}; a symlink
  \texttt{logs/run\_du/latest} is updated each launch.
\end{itemize}
The DU process owns the NGAP connection to the AMF and the GTP-U tunnel to the UPF.

%% --------------------------------------------------------------------------
\subsection{\texttt{run\_ru.sh} --- launch the O-RU (nr-oru)}
%% --------------------------------------------------------------------------

Launches \texttt{openairinterface5g/build/nr-oru} with:
\begin{itemize}[leftmargin=1.4em]
\item Config file \texttt{ru\_test.conf} (FH timing windows \texttt{T2a\_up},
  \texttt{callbacks\_per\_slot=14}, TDD slot counts, IQ compression parameters).
\item CPU affinity to isolated cores 10--15.
\item Extra vrtsim arguments forwarded via \texttt{VRTSIM\_RU\_EXTRA\_ARGS} (used by the
  sweep to pass \texttt{--vrtsim.chanmod}, \texttt{--vrtsim.timescale}, etc.).
\item Log file \texttt{logs/run\_ru/<timestamp>/ru.log}.
\end{itemize}
The RU process owns both DPDK VFs (vf0, vf1 for UL; vf2, vf3 for DL), loads
\texttt{liboran\_fhlib\_5g.so}, and manages the vrtsim shared-memory air interface.

%% --------------------------------------------------------------------------
\subsection{\texttt{run\_ue.sh} --- launch the UE (nr-uesoftmodem)}
%% --------------------------------------------------------------------------

Launches \texttt{openairinterface5g/build/nr-uesoftmodem} with:
\begin{itemize}[leftmargin=1.4em]
\item SIM credentials (IMSI, Ki, OPc) matching the OAI 5GC subscriber database.
\item \texttt{--band 77 --numerology 1 --ssb <offset>} matching the cell configuration.
\item CPU affinity to isolated cores 20--23.
\item Extra vrtsim arguments forwarded via \texttt{VRTSIM\_UE\_EXTRA\_ARGS} (channel model,
  RX-SNR target, etc.).
\item A \texttt{poll\_attach} function that waits for \texttt{oaitun\_ue1} to come up (UE
  has an IP) or times out. The sweep retries this up to \texttt{MAX\_ATTACH\_TRIES} times to
  absorb the $\sim$50\% per-bringup flakiness (Section~2.11).
\end{itemize}

%% --------------------------------------------------------------------------
\subsection{\texttt{kill\_all.sh} --- tear down all OAI processes}
%% --------------------------------------------------------------------------

Sends \texttt{SIGTERM} (then \texttt{SIGKILL} after a timeout) to
\texttt{nr-softmodem}, \texttt{nr-oru}, and \texttt{nr-uesoftmodem}. Does \emph{not}
automatically clear hugepages --- that must be done manually between runs.

%% --------------------------------------------------------------------------
\subsection{\texttt{sweep\_iq\_width.sh} --- main experiment orchestration ($\sim$40\,KB)}
%% --------------------------------------------------------------------------

The sweep script is the primary experiment harness. Its outer loop iterates over a parameter
matrix (iq\_width, compression method, antenna count, channel model, UE speed, FH timing
knobs, bandwidth, RX-SNR target) and for each point:

\begin{enumerate}[leftmargin=1.8em]
\item \textbf{Back up configs.} Copies \texttt{du\_test.conf}, \texttt{ru\_test.conf}, and
  \texttt{ue\_test.conf} to a per-trial log directory and patches them in-place
  (\texttt{min\_grant\_prb}, TDD slots, IQ width, compression header, \texttt{T2a\_up},
  \texttt{T1a\_cp\_ul}, antenna counts, etc.).

\item \textbf{Build channel model (optional).} If \texttt{CHANMOD=1}, the
  \texttt{build\_chanmod()} function generates \texttt{channelmod\_sweep.conf}, inserts an
  \texttt{@include} directive into \texttt{ru\_test.conf} and \texttt{ue\_test.conf}, and
  sets the vrtsim extra-arg variables (\texttt{VRTSIM\_RU\_EXTRA\_ARGS},
  \texttt{VRTSIM\_UE\_EXTRA\_ARGS}).

\item \textbf{Launch DU $\rightarrow$ RU $\rightarrow$ UE} (in that order, with brief waits
  between each). Calls \texttt{run\_du.sh}, \texttt{run\_ru.sh}, \texttt{run\_ue.sh} in
  sub-shells.

\item \textbf{Retry-to-attach loop.} Polls \texttt{ip addr show oaitun\_ue1} until the UE
  has an IP or the attempt count exceeds \texttt{MAX\_ATTACH\_TRIES}. On failure, kills all
  processes, clears hugepages, and retries the full launch from step 3.

\item \textbf{Ensure iperf3 server.} \texttt{ensure\_iperf\_server()} checks whether an
  iperf3 server is already running in the UPF netns; kills stale instances and starts a
  fresh one with \texttt{nsenter}. Only one server may be running at a time.

\item \textbf{Run iperf3 client} on the UE side for the configured duration
  (\texttt{IPERF\_TIME}, default 20~s), UDP (\texttt{IPERF\_UDP=1}),
  offered rate \texttt{IPERF\_UDP\_RATE=20M}.

\item \textbf{Parse metrics} from \texttt{du.log}, \texttt{ru.log}, \texttt{ue.log}, and the
  iperf server log: UL/DL Mbps, jitter, loss, FH rx/tx Mbps, RX SNR
  (\texttt{ULSCH\_STATS}), FH lead-time (\texttt{fh\_lead\_mean}), DL late-drops
  (\texttt{fh\_late\_total}), UL symbol completeness (\texttt{ul\_sym\_present\_pct}).

\item \textbf{Append to \texttt{summary.csv}} and generate \texttt{report.md} for the
  trial directory. The sweep never modifies these files once written.

\item \textbf{Restore configs} from the backup copies, kill all processes, and clear
  hugepages before the next iteration.
\end{enumerate}

Key environment variables controlling the sweep:
\texttt{IQ\_WIDTHS}, \texttt{CHANMOD}, \texttt{CHAN\_TYPE}, \texttt{CHAN\_DS\_TDL},
\texttt{CHAN\_FORGETFACT}, \texttt{NB\_ANT}, \texttt{UE\_COUNT},
\texttt{UE\_SPEED\_KMH}, \texttt{CHAN\_RX\_SNR\_DB}, \texttt{T2A\_MIN\_UP},
\texttt{BW\_PRB}, \texttt{VRTSIM\_TIMESCALE}, \texttt{IPERF\_UDP},
\texttt{IPERF\_UDP\_RATE}, \texttt{IPERF\_MANAGE\_SERVER},
\texttt{MAX\_ATTACH\_TRIES}.

%% ==========================================================================
\section{Experiments \& Results}
%% ==========================================================================

All figures are from \texttt{logs/iq\_width\_sweep/<timestamp>/summary.csv}.
UL is the \textbf{server-side received rate} (iperf3 in the UPF netns as described in
Section~\ref{sec:iperf}), not the client send rate.
The cell is \textbf{24~PRB $\mu$1 = 10~MHz} throughout
(\texttt{carrier\_hz = 4\,049\,760\,000}, band n77).
Outcomes are reported as-is, including the flat and cliff-shaped ones.

%% --------------------------------------------------------------------------
\subsection{Result 0: 10 MHz link works end-to-end \texorpdfstring{\yes}{[OK]}}
%% --------------------------------------------------------------------------

All UL figures use \texttt{IPERF\_UDP=1 IPERF\_UDP\_RATE=20M}; the $\sim$5~Mbps figure is
the radio-limited received rate, well below the 20~Mbit/s offered.
Representative clean 1$\times$1 ideal-channel runs:

\begin{table}[h]\centering\small
\begin{tabular}{@{}l c c c c c c@{}}
\toprule
\textbf{Run} & \textbf{iq} & \textbf{SNR dB} & \textbf{FH tot Mbps}
  & \textbf{UL Mbps} & \textbf{jitter ms} & \textbf{UL loss} \\
\midrule
\texttt{rxsnr\_track\_20dB} & 9  & 20.0 & 150.8 & \textbf{4.995} & 1.83 & 0.00\% \\
\texttt{20260531\_120552}   & 9  & 42.7 & ---   & \textbf{5.004} & ---  & ---    \\
\texttt{20260531\_124630}   & 16 & 45.7 & 254.0 & \textbf{5.006} & ---  & ---    \\
\bottomrule
\end{tabular}
\caption*{\small 1$\times$1, 24 PRB / 10 MHz, ideal channel. UL sym present = 100\% on all clean runs.}
\end{table}

UE syncs, attaches, gets an IP (10.0.0.x), passes NAS, and carries $\sim$5~Mbps UL at 0\%
loss, $\sim$1.8~ms jitter, 100\% UL symbol completeness. Both iq9 (BFP) and iq16
(uncompressed) work.

\textbf{Attach flakiness caveat.} $\approx$50\% of bringups time out (Section~2.11). All
results below are from successful bringups only; the harness retries to get clean data points.

%% --------------------------------------------------------------------------
\subsection{Compression (iq\_width) $\rightarrow$ FH load: linear; UL flat
  \texorpdfstring{\yes/\warn}{[OK/!]}}
%% --------------------------------------------------------------------------

Single sweep \texttt{20260531\_140241} (1$\times$1, ideal channel, all four widths
back-to-back from the same bringup):

\begin{table}[h]\centering\small
\begin{tabular}{@{}c c c c c c c@{}}
\toprule
\textbf{iq} & \textbf{FH rx Mbps} & \textbf{FH total Mbps} & \textbf{SNR dB}
  & \textbf{UL Mbps} & \textbf{FH lead sym} & \textbf{UL loss} \\
\midrule
8  & 94.0  & 135.2 & 38.55 & 4.999 & 12.01 & 0.00\% \\
9  & 104.9 & 150.8 & 38.55 & 4.999 & 12.00 & 0.00\% \\
12 & 137.8 & 197.4 & 38.55 & 5.001 & 12.00 & 0.00\% \\
16 & 180.4 & 256.8 & 38.55 & 5.001 & 12.00 & 0.00\% \\
\bottomrule
\end{tabular}
\end{table}

FH load scales linearly with iq\_width as $(3 \cdot \text{iq} + 1) \cdot \text{num\_prb}$
predicts (iq8$\rightarrow$iq16 nearly doubles the byte rate). \yes\ Clean expected result.

UL throughput, SNR, and FH latency margin are all \textbf{flat}. \warn\ Honest null result:
at 10~MHz / 1$\times$1 the FH is so over-provisioned (180~Mbps into a 6.4~Gbps loopback)
that compression buys nothing on UL. The ``heavy compression helps under FH pressure''
tradeoff only appears when the FH pipe is artificially constrained (Section~4.5) or the load
is driven up with more antennas --- and those hit other walls first (Sections~4.3, 2.12).

%% --------------------------------------------------------------------------
\subsection{Antennas (MIMO) $\rightarrow$ FH scales $\times$N; wide-BW cliffs hard}
%% --------------------------------------------------------------------------

FH load scales with antenna count as expected (iq9, 10~MHz):

\begin{table}[h]\centering\small
\begin{tabular}{@{}c c l@{}}
\toprule
\textbf{Config} & \textbf{FH rx Mbps} & \textbf{Note} \\
\midrule
1$\times$1 & 104.9 & baseline \\
2$\times$2 & 209.9 & 2$\times$ (\texttt{20260530\_221945}: UL 4.991, SNR 42.5) \\
4$\times$4 & 419.8 & 4$\times$ (\texttt{20260530\_211642}) \\
\bottomrule
\end{tabular}
\end{table}

\textbf{2$\times$2 at 10~MHz works} with chanmod ON (mobility):
\texttt{20260531\_144956} $\rightarrow$ UL 4.978, SNR 46.3;
\texttt{20260531\_145140} $\rightarrow$ UL 4.984, SNR 43.7.
Per-run variance is visible: same config at \texttt{20260531\_142504} gave UL 2.67, SNR 32.8
--- bringup-SNR noise, not a config difference.

\textbf{4$\times$4 at 10~MHz can attach} (\texttt{exp\_chanmod\_4x4\_10mhz\_3kmh}: status ok,
UL 4.998, FH 710/1016~Mbps) but is fragile --- earlier symmetric-4$\times$4 attempts gave
UL~=~0 because the server-side DL chanmod (16 convolutions/slot) overran its core even with
the core-pin fix (Section~2.12).

\textbf{Wide-BW is a hard cliff} (see Section~2.10 for root causes):

\begin{table}[h]\centering\small
\begin{tabularx}{\linewidth}{@{}l c c X@{}}
\toprule
\textbf{Target} & \textbf{FH Mbps} & \textbf{Status} & \textbf{Root cause} \\
\midrule
40 MHz (106 PRB)    & 1085 & ue-timeout      & SIB1 NACK / \texttt{prach\_I0=0} \\
20 MHz 4$\times$4   & 637  & ue-timeout      & chanmod overrun \\
20 MHz 4$\times$4   & 2107 & ue-timeout      & chanmod overrun \\
50 MHz 4$\times$4   & 1645 & ue-timeout      & sample-rate wall \\
25 MHz 4$\times$4   & ---  & ru-timeout (crash) & DU realtime \\
\bottomrule
\end{tabularx}
\end{table}

The FH carried up to 2.1~Gbps before the endpoints fell over --- bandwidth is not the limit.

%% --------------------------------------------------------------------------
\subsection{RX SNR $\rightarrow$ UL throughput: clean monotonic tradeoff
  \texorpdfstring{\yes}{[OK]}}
%% --------------------------------------------------------------------------

Using \texttt{--vrtsim.rx-target-snr-db} (time-domain AGC applied on UL only):

\begin{table}[h]\centering\small
\begin{tabular}{@{}c c c c l@{}}
\toprule
\textbf{Target SNR dB} & \textbf{Measured SNR dB} & \textbf{UL Mbps} & \textbf{jitter ms} & \textbf{Run} \\
\midrule
20 & 20.00 & \textbf{4.995} & 1.83 & \texttt{rxsnr\_track\_20dB\_160100} \\
8  & 8.00  & \textbf{2.148} & 4.92 & \texttt{rxsnr\_track\_8dB\_160232}  \\
\bottomrule
\end{tabular}
\end{table}

The AGC hits its target exactly and UL tracks SNR monotonically while jitter rises.
The same shape appears incidentally across chanmod runs:
SNR $\approx$10~dB $\rightarrow$ UL $\approx$1.0--1.7~Mbps;
SNR $\approx$22~dB $\rightarrow$ UL $\approx$2.2~Mbps;
SNR $\geq$38~dB $\rightarrow$ UL $\approx$5.0~Mbps.
This is the cleanest physical tradeoff in the dataset.

%% --------------------------------------------------------------------------
\subsection{FH latency $\rightarrow$ UL: injectable on DL path only; sharp knee then cliff
  \texorpdfstring{\warn}{[!]}}
%% --------------------------------------------------------------------------

The loopback FH is far too fast to create natural latency (Section~2.10), so it was injected.

\textbf{(a) \texttt{T2a\_min\_up} --- the one graceful knob.}
This parameter gates the DL U-plane acceptance window at the O-RU
(\texttt{process\_ru\_uplane} = O-RU receiving DL from the DU). Raising it forces late DL
packets to be dropped; the clean metric is \texttt{up\_late} (RU log ``Packets late: N'').
Sweep: 2$\times$2, 3~km/h, iq9, retry-to-attach:

\begin{table}[h]\centering\small
\begin{tabularx}{\linewidth}{@{}c c c X@{}}
\toprule
\textbf{T2a\_min\_up (\us)} & \textbf{up\_late} & \textbf{UL Mbps} & \textbf{Note} \\
\midrule
200 & 0             & $\sim$4.4 & baseline window \\
300 & 0             & $\sim$5.0 & clean \\
350 & $\sim$8       & $\sim$4.7 & first drops (\texttt{20260601\_004209}: late=8, UL 4.994) \\
400 & $\sim$120     & $\sim$4.0 & knee \\
450 & $\sim$130\,000 & $\sim$2.5 & \textbf{cliff} (retx $\approx$26\%) \\
$\geq$600 & thousands & 0 (no-attach) & DL link breaks (\texttt{20260601\_005323}: late=4000) \\
\bottomrule
\end{tabularx}
\end{table}

Sharp knee at 400$\rightarrow$450~\us\ (drops jump $\sim$1000$\times$). \textbf{This is the
downlink FH path}: UL collapses indirectly because the UE loses DL/DCI. Confirmed by capping
DU egress (\texttt{20260531\_141132}): FH lead-time went 12.0 $\rightarrow$ $-$8.29~sym (DL
arriving after deadline), late=561\,759, UL=0.

\textbf{(b) \texttt{Ta3} / \texttt{Ta4} (true UL FH windows) are no-ops on this loopback.}
Swept Ta4\_max \{440,420,410,405\} and Ta3 \{470,900,1100\}:
\texttt{ul\_sym\_present\_pct} stayed 100\% and UL stayed $\sim$5~Mbps everywhere.
Honest negative result --- they would bite on a real FH link.

\textbf{(c) VF rate-cap is a cliff, not a dial.}
Capping RU/DU egress below the IQ load $\Rightarrow$ mbuf exhaustion $\Rightarrow$ crash
(small pool) or $\sim$74~ms queued latency $\Rightarrow$ no-attach (32$\times$ pool). No
stable graceful regime exists for a CBR fronthaul with a rate limiter.

%% --------------------------------------------------------------------------
\subsection{Mobility $\rightarrow$ UL: barely moves (forgetfact proxy, not real Doppler)
  \texorpdfstring{\warn}{[!]}}
%% --------------------------------------------------------------------------

\texttt{max\_Doppler} is marked ``CURRENTLY NOT IMPLEMENTED'' in \texttt{sim.h:100}. Mobility
is approximated by the \texttt{forgetfact} parameter: 0 = static, 1 = fully decorrelates
each slot. At 3~km/h forgetfact~=~0.0133; at 60~km/h forgetfact~=~0.266.

2$\times$2 TDL-D runs at 3~km/h give UL 4.96--4.98~Mbps when SNR is high
(\texttt{20260531\_144956}/\texttt{145140}). At 60~km/h the result was similarly UL
$\approx$4.99~Mbps. At 120~km/h (forgetfact~=~0.53) bringup fails --- not from UL
degradation but from \textbf{DL initial-sync failure}: the SSB decorrelates faster than
the UE can acquire. There is no significant mobility $\rightarrow$ UL throughput effect at
high SNR on this testbed. UL channel estimation is in-slot (PUSCH DMRS), so there is no
classical CSI-aging mechanism on the UL data path.

%% --------------------------------------------------------------------------
\subsection{Summary}
%% --------------------------------------------------------------------------

\begin{table}[h]\centering\small
\begin{tabularx}{\linewidth}{@{}l l X l@{}}
\toprule
\textbf{Axis} & \textbf{Knob} & \textbf{Effect on UL @ 10 MHz} & \textbf{Verdict} \\
\midrule
Link end-to-end & --- & 5~Mbps, 0\% loss, 100\% UL syms & \yes\ ($\sim$50\% flaky attach) \\
Compression & iq\_width 8$\rightarrow$16 & FH load 94$\rightarrow$180~Mbps; UL flat & \warn\ FH over-provisioned \\
Antennas & 1$\times$1$\rightarrow$4$\times$4 & FH $\times$N; UL $\sim$5; 4$\times$4 fragile & \yes\ FH scales; \warn\ chanmod overrun \\
Bandwidth & $>$24~PRB & --- & \no\ cliff $\rightarrow$ 10~MHz ceiling \\
RX SNR & AGC 20$\rightarrow$8~dB & 5.0 $\rightarrow$ 2.1~Mbps & \yes\ clean monotonic \\
FH latency (DL) & T2a\_min 200$\rightarrow$450~\us & 5.0 $\rightarrow$ 2.5~Mbps, then cliff & \warn\ DL path; sharp knee \\
FH latency (UL) & Ta3 / Ta4 & present 100\%; UL flat & \warn\ no-op on loopback \\
VF cap & egress rate & crash or 74~ms queue & \no\ cliff, not a dial \\
Mobility & forgetfact proxy & flat at high SNR & \warn\ weak; not real Doppler \\
\bottomrule
\end{tabularx}
\caption*{\small The two axes that genuinely bend UL goodput are \textbf{RX SNR} (clean
monotonic) and \textbf{injected DL FH latency via T2a\_min} (sharp knee). All others
correctly change FH load or channel state but do not move UL goodput at 10~MHz because the
binding constraint is the $\sim$5~Mbps grant cap, not the FH pipe nor the channel.}
\end{table}

\appendix

%% ==========================================================================
\section{Git Commits}
%% ==========================================================================

\textbf{Branch \texttt{compression-plus-timing-fix} on \texttt{openairinterface5g}:}

\begin{table}[h]\centering\small
\begin{tabularx}{\linewidth}{@{}l X@{}}
\toprule
\textbf{Commit} & \textbf{Description} \\
\midrule
\texttt{cd509aad8a} & Initial single-machine O-RU recopy working \\
\texttt{057d167ee}/\texttt{dae942c61} & BFP compression; no-comp path works; extra bytes in compression header \\
\texttt{3b10464162} & \textbf{PHY fix} (ULSCH amplitude/timing) --- the FH-tuned PHY that later breaks wide-BW and rfsim SIB1 \\
\texttt{40eaf33c9c} & \textbf{BFP without AVX512} --- scalar reimplementation (Section~2.1) \\
\texttt{565e936c19} & Stability fix (spin-wait, \texttt{>{}>2} scale) \\
\texttt{47dcaef44b} & \textbf{iq16 PUSCH offset} --- skip compression header for \texttt{compMeth=NONE} (Section~2.5) \\
\texttt{37b55fa535} & \textbf{FH UL fix} (reset tti+10) + chanmod core-pin + R-bucket + pthread-setname (Section~2.4) \\
\texttt{a559a1f755} & vrtsim time-domain RX-SNR AGC \\
\bottomrule
\end{tabularx}
\end{table}

\textbf{Branch \texttt{fh-latency-experiment} on \texttt{oran\_lab}:}
\texttt{0e9b835} sync success $\rightarrow$
\texttt{c10a7c0} UL throughput fix $\rightarrow$
\texttt{6a6ad68} compression success $\rightarrow$
\texttt{840e405} stability $\rightarrow$
\texttt{55493cb} reliable UDP UL capture $\rightarrow$
\texttt{fae7cc8} FH-latency harness $\rightarrow$
\texttt{b10743c} RX-SNR AGC wiring.

%% ==========================================================================
\section{Key configuration parameters}
%% ==========================================================================

\begin{table}[h]\centering\small
\begin{tabularx}{\linewidth}{@{}l l l X@{}}
\toprule
\textbf{Parameter} & \textbf{File} & \textbf{Value} & \textbf{Why / issue} \\
\midrule
\texttt{T1a\_cp\_ul}          & du\_test.conf    & \texttt{(285, 535)}          & PRACH C-plane timing (Sec.~2.3a) \\
\texttt{callbacks\_per\_slot} & oran-init.c      & \texttt{14}                  & One RU wakeup per OFDM symbol (Sec.~2.3) \\
Buffer reset offset           & oaioran.c        & \texttt{tti+10}              & MSG3/UL-buffer fix (Sec.~2.4) \\
\texttt{min\_grant\_prb}      & du\_test.conf    & \texttt{16}                  & Break 5-RB grant cap (Sec.~2.7) \\
TDD pattern                   & du/ru\_test.conf & \texttt{D M U U U}           & UL-heavy (Sec.~2.8) \\
\texttt{ul\_prbblacklist}     & du\_test.conf    & \texttt{0,1,2,3,20,21,22,23} & Hide edge PRBs $\rightarrow$ 16 usable (Sec.~2.10) \\
\texttt{tx\_amp\_backoff\_dB} & conf             & \texttt{12}                  & Correct for 7.2 split (RU does IFFT) \\
\texttt{IPERF\_UDP\_RATE}     & sweep env        & \texttt{20M}                 & Offered load; true rate is server-side (Sec.~2.9) \\
\bottomrule
\end{tabularx}
\end{table}

\end{document}
